\chapter{Conclusion}

The long-standing security foundation of RSA – the assumed hardness of factoring – faces a fundamental and qualitative challenge from quantum computing. In 1994, Peter Shor discovered a quantum algorithm that solves the Integer Factorization Problem in polynomial time, specifically \(O((\log n)^3)\) operations. Unlike classical algorithms, Shor's algorithm does not rely on sub-exponential heuristics; it exploits quantum parallelism and the quantum Fourier transform to find the period of a modular exponential function. This period directly reveals the factors of \(n\). If a sufficiently large, fault-tolerant quantum computer were built, Shor's algorithm would break RSA completely, regardless of key size. The security margin provided by increasing the bit-length of \(n\) becomes irrelevant against a quantum adversary, because the running time of Shor's algorithm grows only polynomially with the number of bits. While current quantum computers are far too small and error-prone to factor cryptographically relevant numbers, progress in quantum hardware and error correction is accelerating.

This quantum threat introduces the serious concern of ``store now, decrypt later'': encrypted data captured today can be stored with the expectation that future quantum computers will decrypt it. For information that must remain confidential for decades, this is already a pressing risk. Nevertheless, until large-scale quantum computers become a reality, RSA remains a cornerstone of modern cybersecurity when implemented correctly — using sufficiently large primes, proper padding (e.g., OAEP), and secure key management. In its properly implemented form, RSA is extraordinarily safe for protecting sensitive data today.

We hope this work has been useful and beneficial, and we have endeavoured to cover the subject of RSA cryptosystems and cryptanalysis from all essential angles. As with any technical survey, errors or oversights may remain; if you spot any mistake or wish to suggest improvements, please do not hesitate to inform us. Your feedback is greatly appreciated.